%; whizzy chapter -dvi
% -initex iniptex -latex platex -format platex -bibtex jbibtex -fmt fmt
% 以上 whizzytex を使用する場合の設定。
 
%     Tokyo Debian Meeting resources
%     Copyright (C) 2012 Junichi Uekawa
%     Copyright (C) 2011, 2015 Nobuhiro Iwamatsu

%     This program is free software; you can redistribute it and/or modify
%     it under the terms of the GNU General Public License as published by
%     the Free Software Foundation; either version 2 of the License, or
%     (at your option) any later version.

%     This program is distributed in the hope that it will be useful,
%     but WITHOUT ANY WARRANTY; without even the implied warranty of
%     MERCHANTABILITY or FITNESS FOR A PARTICULAR PURPOSE.  See the
%     GNU General Public License for more details.

%     You should have received a copy of the GNU General Public License
%     along with this program; if not, write to the Free Software
%     Foundation, Inc., 51 Franklin St, Fifth Floor, Boston, MA  02110-1301 USA

%  preview (shell-command (concat "evince " (replace-regexp-in-string "tex$" "pdf"(buffer-file-name)) "&"))

%%ここからヘッダ開始。

\documentclass[mingoth,a4paper]{jsarticle}
\usepackage{monthlyreport}
% 日付を定義する、毎月変わります。
\newcommand{\debmtgyear}{2015}
\newcommand{\debmtgmonth}{10}
\newcommand{\debmtgdate}{18}
% started from zero:
% (let ((year 2013) (month 7)) (+ (* (- year 2005) 12) month -1))
\newcommand{\debmtgnumber}{131}

% tikz picture の為のマクロ設定
\usepackage[dvipdfmx]{graphicx}
\usepackage{tikz}

\begin{document}

\begin{titlepage}
\thispagestyle{empty}
% タイトルページ:編集必要な部分は最初のマクロに飛ばすこと

\vspace*{-2cm}
第\debmtgnumber{}回 東京エリア Debian 勉強会資料\\
\hspace*{-2cm}
\includegraphics{image2012-natsu/dotdeb.pdf}\\
\hfill{}\debmtgyear{}年\debmtgmonth{}月\debmtgdate{}日

% ここはアップデートすること
% 全角文字にしないとフォントのサイズが合わないので注意
\rotatebox{10}{\fontsize{30}{30} {\gt 特集 ： Debconf15 ビデオ紹介}}\\

\vspace*{-2cm}
\hfill{}\includegraphics[height=6cm]{image200502/openlogo-nd.eps}
\end{titlepage}

\newpage

\begin{minipage}[b]{0.2\hsize}
 \definecolor{titleback}{gray}{0.9}
 \colorbox{titleback}{\rotatebox{90}{\fontsize{80}{80} {\gt デビアン勉強会} }}
\end{minipage}
\begin{minipage}[b]{0.8\hsize}
\hrule
\vspace{2mm}
\hrule
\begin{multicols}{2}
\tableofcontents
\end{multicols}
\vspace{2mm}
\hrule
\end{minipage}

\dancersection{事前課題}{野島 貴英}

今回の事前課題は以下です:
\begin{enumerate}
\item 本日、何の作業をやるかを宣言ください。
\item (オプション) どこで今回の勉強会の開催を知りましたか？
\item (オプション) 何について聞きたい／参加者と話をしたいですか？
\end{enumerate}
この課題に対して提出いただいた内容は以下です。
\begin{multicols}{2}
{\small
\begin{prework}{ dictoss }
  \begin{enumerate}
  \item (なし)
  \end{enumerate}
\end{prework}

\begin{prework}{ uwabami }
  \begin{enumerate}
  \item (なし)
  \end{enumerate}
\end{prework}

}
\end{multicols}

\dancersection{最近のDebian関連のミーティング報告}{野島 貴英}

\subsection{第130回東京エリアDebian勉強会 with 第3回 Debianパッケージング道場}

\begin{itemize}
\item 場所は「イベント\&コミュニティスペース dots.」をお借りしての開催でした。
\item 参加者は14名でした。
\item 岩松さんにより、Debianパッケージングの作り方についてのセミナ及びハンズオンが行われました。最後に成果発表をしました。
\end{itemize} 

 Debian公式デベロッパの岩松さんより、昨今のDebianパッケージの開発の作法について説明がありました。文献がなかなか見当たらないgbp(git buildpackage)の仕組みと使い方の説明という、資料としても非常に貴重な発表でした。

 また、東京エリアDebian勉強会としては、初めて会場に「イベント\&コミュニティスペースdots.」を使わせていただきました。イベントの出席登録がしてあれば参加者はその間出入りが自由、コーヒーサーバー併設、無料の無線LANが利用できるという事に加え、内装もモダンな感じであり、勉強会開催にはとても良い場所でした。 

% % (query-replace-regexp "<.*?>" "")
% % (query-replace-regexp "^[	 ]\+" "")

%-------------------------------------------------------------------------------
\dancersection{DebConf15 ビデオ紹介}{野島 貴英}
%-------------------------------------------------------------------------------

\subsection{はじめに}

 毎年1回、世界中のDebian Project関係者及び熱心なユーザらが集まり、ハッカソンをしたり、発表をしたりするイベントとして、DebConfがあります。16回目\footnote{DebConf 0があるため}の開催のDebConf15が、2015/8/15-22の間、ドイツのハイデルベルクで開かれました。

 公式URLは \url{http://debconf15.debconf.org/}となります。

 ここでは、Debconf15で行われたセッションのうち、字幕ファイルが用意されているものについて紹介してみます。

\subsection{DebConf15 セッションのビデオ}

 DebConfでは、Video Teamが各セッションをビデオに撮り公開しているため、いつでもセッションの内容を見ることができます。なお、Debianはフリー（自由）にこだわるため、フリーなフォーマットである、webmが動画フォーマットとして利用されています。

　掲載先：\url{http://debconf15.debconf.org/videostream.xhtml}

 しかしながら、iphone/Androidのスマートフォンで気軽に見たいという今時のニーズもあるかと思います。幸い、youtubeでもDebConf15のビデオが公開されていましたので紹介しておきます。

 youtube:\url{https://www.youtube.com/playlist?list=PLz8ZG1e9MPlz2bUTzfgJhOJCxwT866D4w}
  
\subsection{DebConf15 ビデオ字幕編}

 DebConfは世界中からDebian Project関係者、及び、ユーザが集まるイベントですので、公用語は全て英語になります。発表も英語です。英語を母国語としない人にとってはヒアリングが苦手な方もいらっしゃいます。こういった人のために、現状、数は少ないですが、いくつかの英語の字幕が起こされています。

 字幕取得先:\url{http://ftp.acc.umu.se/pub/debian-meetings/2015/debconf15/subtitles/english/}

 字幕ファイルの使い方は次の通りです。
 
  \begin{description}
\item [Step 1.] 先のURLから、*.srtファイルを取得する。
\item [Step 2.] totem/vlc/mplayerなどでDebConf15の動画を開き、字幕というメニューを選んで対応する.srtファイルを指定します。ファイル名はセッションの名前になっています。
  \end{description}    
  
 字幕付き再生をDebian sid上で行っている様子を載せます。
  
\begin{figure}[H]
\begin{center}
\includegraphics[width=0.5\hsize]{image201510/subtitle.png}
\end{center}
\caption{字幕付き再生例}
\end{figure}

\subsection{今回のビデオ紹介}

 今回紹介予定の具体的なセッション名は、

\begin{itemize}
\item Stretching out for trustworthy reproducible builds
\item Thanks for maintaining a desktop environment. But is it accessible?
\end{itemize}

となります。

\subsection{Stretching out for trustworthy reproducible builds}

 ドイツの有名なChaos Computer Club\footnote{Wikipedia-jpで引いてみて下さい。ドイツの有名なコンピュータ技術のエキスパート集団。}にも所属されているDebian開発者らによる、Reproducible Buildsについてのセッションです。

\begin{figure}[H]
\begin{center}
\includegraphics[width=0.5\hsize]{image201510/reproduct.png}
\end{center}
\caption{Stretching out for trustworthy(略)の発表}
\end{figure}

如何にビデオの内容をかいつまんで紹介します。

\subsubsection{動機}

 \begin{itemize}
 \item The 31st Chaos Communication Congress (31C3)\footnote{Chaos Computer Club主催の毎年行われるイベント}にて、パッケージのバイナリにトロイの木馬が巧妙に仕掛けられているか？を調べるにはReproducible Buildsをしたほうが良いという発表を行ったとのことです。
\item 31C3のわずか数カ月後に今度はEdward Snowdenさんにより、CIAのStrawhourseというコード名に関するCIA conference 2012の内部文章がリークされました。内容は、MacOS/iOSのSDKに不当な改造を行い、生成されるバイナリにCIAが利用するためのトロイの木馬を仕掛けるという驚くべき内容でした。これにより、Reproducible Buildsが益々急務になったとのことです。なお、リーク文章は、\url{https://theintercept.com/document/2015/03/10/strawhorse-attacking-macos-ios-software-development-kit/}で参照できます。
 \end{itemize}

\subsubsection{セキュリティ以外の良い点}

\begin{itemize}
 \item ビルド環境によらず同じバイナリができる、また、クロスビルドの確認ができるようになる、
 \item Debug packageがいつでも（バイナリ作ったあとでも）作れるとか、
 \item FTBFS\footnote{Fails To Build From Sourceの略}が早くわかるとか、
 \item バージョン上げた時の.debの差分が小さくなるとのことです。
 \end{itemize}

\subsubsection{現在のReproducible Builds状況は次の通り}

 \begin{itemize}
 \item Bitcoin/Tor/Corebootは完了している。
 \item Debian/FreeBSD/NetBSD/OpenWrtは進行中。
 \end{itemize}

\subsubsection{工夫と苦労}
 
\begin{itemize}
\item 環境変数SOURCE\_DATE\_EPOCHに時刻(エポック秒)を指定すると、その時刻でビルドしたようにビルドするように様々なツールを改造しupstreamへ提供し取り込んでもらったそうです。なお、これだけでは足らないパッケージが沢山あったらしく、ビルドの日付が埋め込まれる部分をReproducible Builds出来ないとBTSしたりして対策も多数したとのことです。
\item tarにビルド環境の都度のユーザ名、グループ名が混じってしまう件を対策したそうです。
\item ファイルシステムとlocale環境変数(LANG,LC\_ALL変数)との違いによるソートの振る舞いの違い、プログラムの出力が異なってしまう件の対策をしたそうです。
\end{itemize}

言われてみると「なるほど！」と気がつく事ばかりで、考えてみれば相当に苦労するような内容ばかりでした。

\subsubsection{視聴後の所感}

よく、巷では簡単にReproducible Buildsは、パッケージのセキュリティ確認の為と簡単に紹介されますが、実はDebianを構成する重要なソフトウェア・パッケージの多くに手を加えなければ実現できないという大変な偉業を果たしていたという内容でした。思わず、これらの偉業に拍手をしたくなりました。

\subsection{Thanks for maintaining a desktop environment. But is it accessible?}

 Debian ProjectにてAccessibilityを担当されている方の発表となります。
Accessibilityに関しての現状と苦労がわかる発表内容となっています。

 プレゼン資料は：\url{http://brl.thefreecat.org/2015-08-22-debconf.pdf}

\begin{figure}[H]
\begin{center}
\includegraphics[width=0.5\hsize]{image201510/accessibility.png}
\end{center}
\caption{Thanks for maintaining a desktop...(略)の発表}
\end{figure}

\subsubsection{Accessibilityについてうっかりすると忘れがちになる大変重要な事}

　Free Softwareは、問題があったり、気に入らなかったら、自分で直せるということが基本であるが、Accessibilityの機能を必要としている人は、基本治したくても直せない場合が多いので、コミュニティーによる修正・改善が必須です。

\subsubsection{Linux Desktop環境のAccessibilityの現状}

\begin{itemize}
  \item Linuxで動作するDesktop環境は、GNOMEがAccessibilityが最もよくできている状況です。
  \item しかしながら、GNOME3を持ってしても、Windowsに比べると10年単位で遅れており、Appleの製品に比べると石器時代の代物と言われても仕方が無い状況です。
  \item 弱視の人には、合成音声によるサポートは厳しい場合（そもそも発音しにくいワードの場合など）があるため、理想的には、Piezo braille cellをサポートすべきです。
\end{itemize}

\subsubsection{Linux Desktop環境のAccessibilityの現状の仕組みと開発方法}

先のプレゼン資料、及び、ビデオにて

\begin{itemize}
  \item LinuxのAccessibilityのフレームワーク、
  \item LinuxのAccessibilityのテスト環境など
 \end{itemize}

が紹介されています。先述のプレゼン資料を参照ください。Accessibilityがどのようにできていて、どうテストすべきかについて、非常に良い資料となっています。

\subsection{おわりに}

　紙面の関係で今回はセッションを２つのみ紹介しましたが、他にも非常に興味深いセッションがあります。なにぶん英語のセッションなので、見て理解するのに苦労する状況ですが、その努力を払っただけの収穫があるのが、Debconfのビデオです。是非、１つ見て、内容を考えてみませんか？きっと、今まで見えていた世界ががらっと変わるような体験が出来ると思っています。

%-------------------------------------------------------------------------------
\dancersection{会場での無線LANのつなぎ方}{野島 貴英,Roger}
%-------------------------------------------------------------------------------
 \subsection{はじめに}

　今回試験として、会場側でフィルタ無しのグローバル回線を用意しました。
ただ、会場側のセキュリティポリシーにより、wpa-psk AES hidden SSIDという
方式での提供となります。

　以下にDebianマシンでの接続方法を記載します。

 また、自分の環境では違うやり方でつながったという方は、野島まで
教えて下さい。こちらでもノウハウとして溜めていく予定です。

 \subsection{wpasupplicant及び/etc/network/interfacesを利用の場合}

 もっとも良いマニュアルは、/usr/share/doc/wpasupplicant/README.Debian.gz
となります。困った場合はこちらも合わせてご参照下さい。

　以下に/etc/network/interfacesの定義について会場の例を記載します。

\begin{commandline}
$ sudo vi /etc/network/interfaces
-----以下のエントリがなければ追記ここから----------
iface wlan0_debian inet dhcp
     wpa-conf /etc/wpa_supplicant/wpa_supplicant_debian.conf
-----以下のエントリがなければ追記ここまで----------
$ sudo vi /etc/wpa_supplicant/wpa_supplicant_debian.conf
-----以下のエントリを追記ここから----------
network={
     ssid=<<会場のSSID>>
     psk=<<会場のパスワード>>
     scan_ssid=1
}
-----以下のエントリを追記ここまで----------
$ sudo chmod 600 /etc/wpa_supplicant/wpa_supplicant_debian.conf
$ sudo ifup wlan0=wlan0_debian
\end{commandline}
%$

 また、ハマってしまった時のデバッグ方法は、
/usr/share/doc/wpasupplicant/README.Debian.gz中の''4. Trubleshooting''の章が便利です。

 \subsection{その他の無線LAN用パッケージを利用の場合}

　すみません、自分が情報を持たないため、現場で教えて下さい。

\cleartooddpage

\vspace*{15cm}
\hrule
\vspace{2mm}
\includegraphics[width=2cm]{image200502/openlogo-nd.eps}
\noindent \Large \bf Debian 勉強会資料\\
\noindent \normalfont \debmtgyear{}年\debmtgmonth{}月\debmtgdate{}日 \hspace{5mm}  初版第1刷発行\\
\noindent \normalfont 東京エリア Debian 勉強会 （編集・印刷・発行）\\
\hrule

\end{document}
